\documentclass[11pt]{article}
\usepackage{amsmath, amssymb, amsthm}
\usepackage{geometry}
\geometry{margin=1in}
\title{CSE 400: Fundamentals of Probability in Computing\Lecture 21-22 Notes (From Chebyshev's Inequality)}
\author{Scribe Refactoring}
\date{}

\begin{document}
\maketitle

\section{Chebyshev's Inequality}

\textbf{Definition:} For any random variable $X$ with mean $\mu$ and variance $\sigma^2$,
[
\Pr(|X - \mu| \geq k\sigma) \leq \frac{1}{k^2}, \quad k > 0
]

\textbf{Equivalent Form:}
[
\Pr(|X - \mu| \geq t) \leq \frac{\sigma^2}{t^2}
]

\subsection{Proof}
Let $Y = (X - \mu)^2$. Using Markov's inequality:
[
\Pr((X-\mu)^2 \geq t^2) \leq \frac{\mathbb{E}[(X-\mu)^2]}{t^2} = \frac{\sigma^2}{t^2}
]

\section{Interpretation}
Chebyshev's inequality provides a universal bound that applies to \textit{any} distribution with finite variance.

\section{Example}
If $\sigma = 2$, then:
[
\Pr(|X - \mu| \geq 4) \leq \frac{1}{4}
]

\section{Chernoff Bound}

\textbf{Motivation:} Provides exponentially tighter bounds than Chebyshev.

\textbf{Setup:} Let $X = \sum X_i$ where $X_i$ are independent Bernoulli.

\textbf{Upper Tail:}
[
\Pr(X \geq (1+\delta)\mu) \leq e^{-\frac{\mu \delta^2}{3}}, \quad 0 < \delta < 1
]

\textbf{Lower Tail:}
[
\Pr(X \leq (1-\delta)\mu) \leq e^{-\frac{\mu \delta^2}{2}}
]

\section{Proof Sketch}
Uses moment generating function (MGF) and Markov's inequality:
[
\Pr(X \geq a) = \Pr(e^{tX} \geq e^{ta}) \leq \frac{\mathbb{E}[e^{tX}]}{e^{ta}}
]

Optimize over $t$.

\section{Poisson Tail Bounds}
For $X \sim \text{Poisson}(\lambda)$:
[
\Pr(X \geq (1+\delta)\lambda) \leq \left(\frac{e^{\delta}}{(1+\delta)^{1+\delta}}\right)^{\lambda}
]

\section{Jensen's Inequality}

\textbf{Definition:} If $f$ is convex, then:
[
f(\mathbb{E}[X]) \leq \mathbb{E}[f(X)]
]

\subsection{Convex Function}
$f$ is convex if:
[
f(\theta x + (1-\theta)y) \leq \theta f(x) + (1-\theta)f(y)
]

\section{Applications}
\begin{itemize}
\item Bounding tail probabilities
\item Performance guarantees in randomized algorithms
\item System reliability analysis
\end{itemize}

\section{Case Study: Server Load}
Tail bounds help estimate worst-case demand in systems.

\section{Summary for Exam}
\begin{itemize}
\item Chebyshev: Works for all distributions
\item Chernoff: Strong exponential bounds
\item Jensen: Links convexity and expectation
\item Tail bounds: Essential for worst-case guarantees
\end{itemize}

\end{document}
